% Generated only from audited measured results.
\newcommand{\OnlineRuns}{840}
\newcommand{\TotalProcesses}{2,712}
\newcommand{\TotalRows}{477,096}
\newcommand{\UsefulInstances}{420,000}
\newcommand{\BestFixedArm}{SS}
\newcommand{\PriorWall}{35.97}
\newcommand{\PriorCPU}{83.89}
\newcommand{\AbstractResults}{Four-candidate screening saves 65.0\% of tuning wall time and 65.9\% of tuning CPU time, with a geometric mean payload penalty of 10.2\% to the same-block empirical best. Its additional payload loss relative to full search's chosen schedule is -1.1\% (95\% CI -3.3\%--0.9\%). For 50 useful executions and an available prior, its total-cost ratio is 1.379 relative to full search (95\% CI 1.329--1.435), where ratios above one favor screening. With calibration counted as useful work, the four-candidate wall ratios are 0.970 at 50 executions and 1.000 at 500, limiting the savings claim to the extra-calibration operating mode. }
\newcommand{\SelectionResults}{Guided 2, Guided 4, and Guided 8 save 82.4\%, 65.0\%, and 31.9\% of tuning wall time, respectively. Their corresponding CPU savings are 83.2\%, 65.9\%, and 32.6\%. Prior-only selection saves 91.1\% wall time and 92.5\% CPU time. These are paired geometric cost ratios over complete tuning phases, not ratios of median milliseconds.}
\newcommand{\QualityResults}{Guided 4's geometric mean payload loss is 10.2\% (95\% CI 7.5\%--13.3\%). It retains 90.7\% of reference speed. Relative to the schedule actually selected by Full 12, its extra payload loss is -1.1\% (95\% CI -3.3\%--0.9\%). Its payload is within 5\% in 44.2\% of runs and within 10\% in 64.2\%. Full search itself has a geometric residual of 11.4\%, illustrating the gap between a finite noisy sweep and an independent best-arm reference.}
\newcommand{\DeploymentResults}{Guided 4's wall-cost ratios at horizons 10, 50, 100, and 500 are 2.014, 1.379, 1.200, 1.053. At 50 executions the CPU-cost ratio is 1.384 (95\% CI 1.335--1.440); at 500 it is 1.053. Ratios above one indicate lower total cost for screening.}
\newcommand{\WorthResults}{At 50 useful executions, 44.2\% of Guided 4 runs satisfy (95\% cluster-bootstrap interval 27.5\%--60.8\%) all three conditions of lower wall cost, lower CPU cost, and at most 5\% payload loss; with a 10\% tolerance this rises to 64.2\%. At 500 executions the corresponding rates are 42.5\% and 60.0\%.}
\newcommand{\CrossoverResults}{For Guided 4, 64 of 120 paired runs have both a tuning saving and a positive payload mean gap to full search. Their median estimated finite crossover is 935 useful executions. This summarizes only runs with a positive gap; it is not a universal operating cutoff.}
\newcommand{\AmortizationResults}{For Guided 4 at horizon 50, the estimated break-even is 488 deployments for wall time and 409 for CPU time. At horizon 500 the estimates are 531 and 448, respectively. When both resources matter, the larger finite count is the relevant lower bound.}
\newcommand{\SensitivityResults}{Under the cross-fitted reference, Guided 4's geometric payload loss becomes -1.7\%; its within-5\% and within-10\% rates become 85.0\% and 92.5\%. The difference from the primary reference quantifies sensitivity to choosing a minimum among noisy same-block timings.}
\newcommand{\ControlResults}{Random 4 has a geometric payload loss of 11.2\% and a horizon-50 wall-cost ratio of 1.403. The trained fixed arm has a loss of 10.2\% and ratio 1.704, with no online calibration. Thus the value of prediction must be considered alongside the simpler ways to spend less on search.}
\newcommand{\DecisionResults}{For the four-candidate policy, the measured exchange is 65.0\% less tuning wall time and 65.9\% less tuning CPU time for 10.2\% aggregate residual payload loss. Its horizon-50 total wall ratio is 1.379, whereas its horizon-500 ratio is 1.053. These measurements make the trade-off explicit; the tolerance hit rate and prior cost determine whether it is acceptable for a particular deployment.}
\newcommand{\ConclusionResults}{On the measured four-core system, Guided 4 avoids two thirds of calibration evaluations and saves 65.0\% tuning wall time and 65.9\% tuning CPU time. Its geometric payload penalty is 10.2\%. For an available prior, its total wall-cost ratio is 1.379 at 50 useful executions and 1.053 at 500. Prior acquisition and tail losses remain part of the decision rather than being hidden costs. The separate equal-useful-work ratios of 0.970 and 1.000 do not establish an in-situ benefit. }
\newcommand{\InsituRuns}{432}
\newcommand{\InsituInstances}{118,800}
\newcommand{\InsituResults}{For 50 equal useful executions, Guided 4's wall-cost ratio is 0.970 (95\% CI 0.933--0.999) and its CPU-cost ratio is 0.970 (95\% CI 0.932--1.001). For 500 equal useful executions, Guided 4's wall-cost ratio is 1.000 (95\% CI 0.987--1.013) and its CPU-cost ratio is 1.001 (95\% CI 0.987--1.014).}
\newcommand{\InsituDecision}{the four-candidate wall-cost ratios are 0.970 at 50 useful calls and 1.000 at 500. These do not establish a time/CPU saving from reduced search in that mode. Avoided candidate tests cannot be counted as avoided necessary executions; the primary extra-calibration savings apply only when calibration is additional work.}
